% results.tex -- input by main.tex
\section*{Results}

\subsection*{Isolation methods pair inoculation with a serum step-down}
In the plotted dual-FBS corpus ($n=169$ of 170 rows parsed), pre-inoculation FBS had a mean of 9.30\% and median 10.0\%; post-inoculation FBS had a mean of 2.24\% and median 2.0\%. Pre-inoculation values were 10\% in 139 of 169 papers; post-inoculation values were 2\% in 113 of 169, with a tail including 0\%, 1\%, 5\% and 10\%. The eight institutional recipes cluster at the same split (about 10\% growth to about 2\% maintenance; one measles recipe at 7.5\% to 2\%) (Fig~1). These frequencies describe practice as sampled and do not rank recipes.


\textbf{Fig 1.} Pre- and post-inoculation FBS in the dual-FBS corpus with institutional recipe overlay.

\subsection*{Guidelines prescribe the same step-down for different reasons}
The five guidelines that prescribe a serum step-down at inoculation agree on the target (about 2\%) but not on the reason (Table~\ref{tab:rationale}). ATCC attributes it to serum interfering with viral attachment; CLSI to cell viability, slower overgrowth and serum-borne viral inhibitors; ASM, NEADL and WOAH give no reason. None of the quoted passages cites an experiment showing that the step-down leaves the morphology of uninoculated cells unchanged.

% rationale.tex -- Table: stated rationale for the serum step-down, by guideline. Input by sections/results.tex.
\begin{table}[!ht]
\centering
\caption{Serum step-down prescribed at or near inoculation, and the rationale each guideline gives for it. Quotations are verbatim.}
\label{tab:rationale}
\begin{tabular}{p{0.2\linewidth}p{0.17\linewidth}p{0.53\linewidth}}
\hline
Guideline & Serum change & Stated rationale \\
\hline
ATCC Virology Culture Guide & to 2\% for inoculation & ``ATCC recommends decreasing the percentage of serum to 2\% as it can interfere with viral attachment.'' \\
CLSI M41-A & about 10\% to 2\% (1--3\%) & ``Medium with the lower concentration of FBS will maintain cell viability and will slow monolayer overgrowth.'' ``Serum may also contain specific and nonspecific viral inhibitors and their presence may vary by lot.'' ``The lower serum concentration supports cell viability while permitting virus replication.'' \\
ASM CPE protocols & growth medium to 2\% & None stated. \\
NEADL PPRV SOP & 10\% to 2\% & None stated. \\
WOAH Terrestrial Manual (2021), ch.\ 3.5.2 & to 2\% & None stated. \\
\hline
\end{tabular}
\end{table}


\subsection*{Matched-medium negative controls are rarely reported}
Of 158 Isolation papers, 49 (31.0\%) reported an Isolation-stage negative control (Tiers A+B). Only 11 (7.0\%) reported one handled in the matched maintenance medium (Tier A); 38 reported an Isolation-stage control without a stated medium match (Tier B), 14 reported controls only at another stage (Tier C) and 95 (60.1\%) reported no qualifying control (Tier D). None of the 11 Tier A papers reported CPE in its negative control; 8 stated the control was CPE-negative and 3 did not say. Absence or mismatch of controls is a methods-reporting finding; it does not by itself establish that CPE is non-specific in any particular laboratory system.

\subsection*{Non-isolation protocols and practice}
USDA APHIS VIRPRO1013, a master-seed extraneous-agent testing protocol, prescribes growth medium with about 5--10\% serum and a later medium with 5--20\% serum, so serum is held or raised rather than stepped down as in Isolation recipes. \pending{Expanded protocol search (R4) and practice reports to follow.}
\assumed{Protocols released under the public-records request show reference-laboratory practice centred on a shift from 10\% growth serum to a mean of about 15\% FBS in the culture medium used at or after inoculation. The distribution is broad: many protocols hold serum at 10\% (10\%$\rightarrow$10\%), many raise it to 20\% (10\%$\rightarrow$20\%), and few or none lower it to 5\% (10\%$\rightarrow$5\%). If this holds, reference-laboratory practice departs from the published 10\%$\rightarrow$2\% step-down; Isolation readouts in those laboratories would then not carry the low-serum confound measured in R3, and the R3 result applies to the published-recipe setting rather than to all practice. Replace with the coded records: number of protocols, distribution of pre/post values, and any stated rationale.}

\subsection*{Lowering FBS alone produces CPE-classified morphology in uninoculated Vero E6 cells}
\emph{Contract laboratory descriptions.} On the 22 frames with their own CRO description, CPE-type terms were named on 1 of 9 Culture A (10\% FBS) frames and 10 of 13 Culture B (2\% FBS) frames. In Culture B the terms were rounded (8 frames), dying cells (7), granularity (6), vacuolation (4), detached (2) and refractile (1); the single Culture A mention was vacuolation. The mean share of the six CPE-type terms per frame was 0.02 in Culture A and 0.36 in Culture B.

\emph{Blinded raters, structured checklist (Pass 2).} Both raters scored more CPE-type morphology on Culture B. The mean CPE-item score sum (maximum 7) was 0.87 (A) vs 1.65 (B) for IR1 and 2.41 (A) vs 4.91 (B) for IR2. Items with the largest B$-$A incidence gaps were ballooned/enlarged (IR1 +0.62; IR2 +0.78) and cytoplasmic strands (IR1 +0.82); for IR2 also nonspecific degeneration (+1.00), syncytia (+0.55) and cell death (+0.22). IR2 rated 20 of 49 Culture B frames as not healthy overall and 29 as partially healthy, against 50 of 50 Culture A frames healthy; IR1 rated all frames healthy overall. IR2's Culture B CPE score rose monotonically from day 1 (3.48) to day 5 (6.34) with Culture A flat; IR1 showed a constant B$>$A offset across days.

\emph{Blinded raters, free text (Pass 1).} With the shared lexicon, IR2 named a CPE-type term on 19 of 50 Culture A and 25 of 49 Culture B frames; IR1's free text named rounded and refractile cells on every frame of both cultures, so its Pass-1 contrast is saturated.

\emph{Briefing-condition pilot arm (IR3v).} \prov{Under the clarified coding, a term from the CRO six was named on 36 of 50 Culture A and 40 of 49 Culture B frames: vacuolation 31 vs 30, rounded/shrinking 8 vs 12, refractile 17 vs 21, detached 0 vs 2 (shared lexicon, any CPE term: 38 vs 42). The single-large-vacuole observation, coded separately, was more frequent in Culture B (14 vs 24 frames), as were uncoded mentions of projections (10 vs 32). Culture B led on days 1--2 (7 and 8 of 10 frames vs 2 and 4) and both cultures reached 8--10 frames per day from day 3. On the 22 CRO-described frames IR3v named a CRO-six term on 5 of 9 (A) and 8 of 13 (B). Direction matches the other sources (B $\geq$ A); the gap (+10 percentage points) is similar to IR2 Pass 1 and smaller than for the CRO. Provisional until IR3v Pass 2.} \pending{IR3v Pass 2 (in progress).}

\textbf{Fig 2.} \todoA{Select the R3 figure set (candidates c0, c2, g2, h, j in \texttt{r3\_image\_labeling/wrk/viz\_figures\_review.md}).}

Table~\ref{tab:r3-cpe-by-source} summarises the per-culture contrast for every source.

% r3-cpe-by-source.tex -- Table: CPE-type morphology by source and culture (R3). Input by sections/results.tex.
\begin{table}[htbp]
\centering
\small
\caption{CPE-type morphology named or marked per frame, by source and culture (Culture A = 10\% FBS, Culture B = 2\% FBS; no inoculum in either). Sources: CRO report and rater label tables (see Data availability).}
\label{tab:r3-cpe-by-source}
\begin{tabular}{llll}
\toprule
\textbf{Source} & \textbf{Measure} & \textbf{Culture A} & \textbf{Culture B} \\
\midrule
CRO & frames naming any of the CRO six (own description) & 1/9 & 10/13 \\
IR1 & Pass 2 mean CPE-item score sum (max 7) & 0.87 & 1.65 \\
IR2 & Pass 2 mean CPE-item score sum (max 7) & 2.41 & 4.91 \\
IR1 & Pass 1 frames with any CPE term (shared lexicon) & 50/50 & 49/49 \\
IR2 & Pass 1 frames with any CPE term (shared lexicon) & 19/50 & 25/49 \\
IR3v & Pass 1 frames with any CRO-six term (clarified)\textsuperscript{a} & 36/50 & 40/49 \\
IR3v & Pass 2 & \multicolumn{2}{l}{in progress} \\
\bottomrule
\end{tabular}
\\[2pt]{\footnotesize \textsuperscript{a} Provisional (briefing-condition pilot arm; Pass 2 in progress).}
\end{table}


\subsection*{Inter-rater agreement}
Raters agreed on direction (B $\geq$ A on aggregate CPE scores) but not on which features carry the difference: IR1 used vacuolation on no frame while IR2 marked it on every frame; IR2 marked syncytia and nonspecific degeneration on Culture B while IR1 marked neither. Feature-identity agreement is not required for the claim tested here, which concerns the direction of CPE-classified morphology between cultures that differ only in serum. Per-item agreement is summarised in Table~\ref{tab:agreement}.

% agreement.tex -- Table: IR1 vs IR2 per-item agreement on the Pass 2 checklist (99 frames). Input by sections/results.tex.
\begin{table}[!ht]
\centering
\caption{Agreement between blinded raters IR1 and IR2 on the Pass 2 checklist (99 frames). Positive = any mark other than No. $\kappa$ is Cohen's kappa on present/absent; weighted $\kappa$ uses three levels (No; Mild/Partial; Yes) with linear weights. A dash means $\kappa$ is undefined because at least one rater gave the same mark on every frame.}
\label{tab:agreement}
\begin{tabular}{lrrrrr}
\hline
Item & IR1 pos. & IR2 pos. & Agree (\%) & $\kappa$ & weighted $\kappa$ \\
\hline
Healthy overall & 99 & 79 & 80 & -- & -- \\
Nuclei defined & 99 & 99 & 100 & -- & -- \\
Cytoplasmic extensions & 99 & 98 & 99 & -- & -- \\
Vacuolation & 0 & 99 & 0 & -- & -- \\
Granularity & 99 & 99 & 100 & -- & -- \\
Ballooned/enlarged & 68 & 38 & 70 & 0.442 & 0.346 \\
Syncytia & 0 & 27 & 73 & -- & -- \\
Cytoplasmic strands & 54 & 99 & 54 & -- & -- \\
Cell death & 3 & 88 & 14 & 0.008 & 0.005 \\
Nonspecific degeneration & 0 & 49 & 50 & -- & -- \\
\hline
\end{tabular}
\end{table}

